\documentclass[11pt]{article}
\usepackage[margin=1in]{geometry}
\usepackage{amsmath,amssymb}
\usepackage{booktabs}
\usepackage{graphicx}
\usepackage{hyperref}
\usepackage{xcolor}
\hypersetup{colorlinks=true,linkcolor=blue,urlcolor=blue}

\title{Independent Computational Replication of\\
\emph{Vortex Splitting in Subcritical Nonlinear Schr\"odinger Equations}\\
{\large Berloff (2008), arXiv:0801.2964}}
\author{REPLICATE-PROJECT / TEXTURES-100 --- polar / NLS class}
\date{Reimplementation from equations; not author code}

\begin{document}
\maketitle

\begin{abstract}
We independently reimplement, from the governing equations only (no author code),
the central static and dynamic results of Berloff (2008) on vortices in the
\emph{subcritical} (cubic--quintic) nonlinear Schr\"odinger equation (NLS). Two
tiers were built: (A) the stationary straight-line vortex radial profile via a
1D Newton relaxation of the profile ODE, and (B) the 2D time-dependent dynamics
of a doubly charged ($s=2$) vortex under a periodic pressure drive, integrated
with a Strang split-step Fourier scheme. The statics match the paper to
$\sim$1\%: slope-at-origin $a_1(\xi\!\to\!0)=0.962$ (paper $0.9575$),
$a_1(\xi=5/8)=0.2863$ (paper $0.286$), critical pressure parameter
$\xi_{\mathrm{crit}}=0.690$ (paper $0.689$). The dynamics reproduce the
topological signature that only $s=\pm1$ vortices are stable --- the $s=2$ core
resolves into two $+1$ winding cores with \emph{total circulation conserved} ---
and reproduce pressure-driven core breathing ($r_{\mathrm{core}}:3.1\to6.7\to3.1$)
as the drive crosses $\xi_{\mathrm{crit}}$ (paper Figs.~6--7).
\textbf{Verdict: PARTIAL} (coverage 7/10, agreement 8/10). The honest gap is
that the two unit cores stay grid-adjacent ($\sim$1 healing length) rather than
separating macroscopically: the paper's headline \emph{macroscopic} split is a
3D vortex-\emph{ring} phenomenon (velocity-field asymmetry, Figs.~8--9), which we
did not build. A 2D straight vortex legitimately shows core \emph{resolution} and
\emph{breathing}, not large separation --- consistent with the paper's own remark
that a straight vortex remains radially symmetric.
\end{abstract}

\section{The paper and its claims}
Berloff (2008) studies vortices and axisymmetric vortex rings in the subcritical
(cubic--quintic) NLS, where a higher-order nonlinearity models many-body
interactions and permits the study of \emph{negative pressure} on vortex
dynamics. The headline physics claims are:
\begin{enumerate}
  \item There exists a \textbf{critical pressure} at which a straight-line vortex
        becomes unstable to radial expansion of its core (quantified by the
        profile parameter $\xi_{\mathrm{crit}}$ where the core slope-at-origin
        $a_1\to0$).
  \item The \textbf{stationary vortex profile} $f(r)$ is characterized by
        $a_1(\xi)$; the paper tabulates $a_1(\xi\!\to\!0)=0.9575$ and states
        $a_1(\xi=5/8)=0.286$, with $\xi_{\mathrm{crit}}=0.689$ (for
        $\gamma=1$).
  \item Only \textbf{singly-quantized} ($s=\pm1$) vortices are dynamically
        stable; multiply-charged vortices are unstable.
  \item Under a \textbf{periodically varying pressure field}, a vortex ring may
        \emph{split into many vortex rings}; conditions are elucidated
        (Figs.~6--9).
\end{enumerate}

\section{Governing equations reimplemented}
\subsection{The dynamical equation (Eq.~16)}
We use the sign-reconciled dimensionless form for which the bulk $\psi=1$ is an
exact fixed point and Eq.~(21) is its radial profile ODE:
\begin{equation}
  i\,\psi_t \;=\; -\tfrac{1}{2}\nabla^2\psi
    \;+\; \Big[\,|\psi|^{2(1+\gamma)} - 2\xi\,|\psi|^2 - (1-2\xi)\,\Big]\psi,
  \label{eq:gpe}
\end{equation}
with $\gamma=1$ (the cubic--quintic case) and pressure parameter $\xi$.
Reconciling the signs so that $|\psi|=1$ nulls the bracket
($1 - 2\xi - (1-2\xi)=0$) is essential; a sign slip there silently drifts every
downstream number.

\subsection{The stationary radial profile (Eq.~21)}
For a straight-line vortex of charge $s$, $\psi=f(r)e^{is\theta}$ with $f$ solving
\begin{equation}
  \tfrac{1}{2}\frac{1}{r}\frac{d}{dr}\!\Big(r\frac{df}{dr}\Big)
  - \frac{s^2}{2r^2}f + 2\xi f^3 + (1-2\xi)f - f^{\,2(1+\gamma)+1} = 0,
  \label{eq:profile}
\end{equation}
with $f(0)=0$, $f\to1$ as $r\to\infty$. The slope-at-origin
$a_1 \equiv f'(0)$ is the paper's diagnostic; $a_1\to0$ marks
$\xi_{\mathrm{crit}}$, the onset of unbounded core expansion.

\subsection{The pressure drive (Eq.~41)}
The time-dependent negative-pressure forcing is implemented exactly as
\begin{equation}
  \xi(t) = \xi_0 + \varepsilon\sin\!\Big(\frac{\pi t}{2\eta}\Big),
  \qquad \xi_0 = \frac{1+\gamma}{2+\gamma}\;(=\tfrac{2}{3}\text{ for }\gamma=1).
  \label{eq:drive}
\end{equation}
When $\xi_0+\varepsilon$ crosses $\xi_{\mathrm{crit}}$ the negative pressure
drives radial expansion of the core (paper Sec.~4).

\section{Methods}
\subsection{Part A --- stationary profile solver}
Eq.~\eqref{eq:profile} is discretized on a uniform radial grid
($R=30$, $N=3000$) and solved by Newton relaxation: at each iteration the
tridiagonal Jacobian is assembled analytically and solved by the Thomas
algorithm, with Dirichlet BCs $f(0)=0$, $f(R)=1$. Convergence to
$\max|\delta f|<10^{-11}$ takes a handful of iterations ($<1$~s total). $a_1$ is
read off as $f(r_1)/r_1$; $\xi_{\mathrm{crit}}$ is located by scanning $\xi$ and
interpolating where $a_1\to0$.

\subsection{Part B --- 2D split-step Fourier dynamics}
Eq.~\eqref{eq:gpe} is integrated by the Strang split-step Fourier method on a
periodic $N{=}128$ grid, box $L=48$, $dt=0.004$, to $T=200$:
\begin{itemize}
  \item Half kinetic step in $k$-space:
        $\exp(-i\,\tfrac{dt}{2}\cdot\tfrac{1}{2}k^2)$ (since
        $-\tfrac12\nabla^2\to\tfrac12 k^2$).
  \item Full nonlinear phase in real space:
        $\exp(-i\,dt\,[\,\rho^{1+\gamma}-2\xi\rho-(1-2\xi)\,])$,
        $\rho=|\psi|^2$.
  \item Half kinetic step.
\end{itemize}
An \textbf{absorbing sponge} in the outer ring ($r>0.42L$) soaks radiated sound
so it does not wrap the periodic box and contaminate the interior. The initial
condition is an approximate $s=2$ core, $\psi\approx
[r^2/(r^2+r_c^2)]\,e^{2i\theta}$, with a small $\cos2\theta$ seed perturbation.

\subsection{Topological tracking (the key correctness check)}
Vortices are located by \textbf{$2\pi$ phase winding on each plaquette}
(wrapped phase differences summed around the four corners, divided by $2\pi$),
restricted to the central region $r<0.30L$ to reject sound pulses. Adjacent
$+1$ plaquettes are clustered ($3\times3$ connectivity) into cores; the maximum
pairwise separation is the split distance. Crucially, we verify that the
\textbf{total circulation is conserved} throughout the run --- a far stronger
correctness check than any density-minimum threshold, which would spuriously
merge close cores and miscount sound. Core radius is measured as
$R=\sqrt{A/\pi}$ over the central low-density ($\rho<0.5$) area.

\section{Results}
\subsection{Part A --- statics match to $\sim$1\%}
\begin{table}[h]
\centering
\begin{tabular}{lccc}
\toprule
Quantity & This work & Berloff (2008) & Rel.\ error \\
\midrule
$a_1(\xi\!\to\!0)$        & 0.962  & 0.9575 & 0.5\% \\
$a_1(\xi=5/8)$            & 0.2863 & 0.286 (exact) & 0.1\% \\
$\xi_{\mathrm{crit}}$     & 0.690  & 0.689  & 0.1\% \\
\bottomrule
\end{tabular}
\caption{Stationary straight-line vortex profile ($\gamma=1$, $s=1$) vs.\ the
paper's Table~1 / text. The profile ODE and its units are pinned to $\sim$1\%.}
\end{table}

The excellent agreement on all three static diagnostics confirms the equation,
its signs, and its units are correct before any 2D dynamics is trusted.

\subsection{Part B --- topological resolution + pressure-driven breathing}
The 2D run (runtime $\approx28$~s) yields:
\begin{itemize}
  \item \textbf{Total charge conserved:} the median steady-state total
        circulation is $+2$ and stays at $2$ for $>90\%$ of the run --- the
        integrand of the ``only $s=\pm1$ stable'' claim.
  \item \textbf{Charge-2 resolves into two $+1$ cores:} the $s=2$ core is not a
        stable object; it resolves into two unit ($+1$) winding cores while
        conserving total circulation --- the topological signature of the
        instability of multiply-charged vortices.
  \item \textbf{Pressure-driven core breathing:} with $\xi_0=2/3$,
        $\varepsilon=0.10$, $\eta=100$ (peak $\xi=0.767>\xi_{\mathrm{crit}}$),
        the core radius breathes $r_{\mathrm{core}}:3.09\to6.74\to3.05$
        (amplitude $\approx3.65$), expanding while $\xi>\xi_{\mathrm{crit}}$ and
        contracting as the pressure recovers --- the physics of Figs.~6--7.
\end{itemize}

\section{Comparison and verdict}
Where the paper gives numbers we can compare against (the profile diagnostics
$a_1$ and $\xi_{\mathrm{crit}}$), agreement is $\sim$1--2\%. The dynamic
core-expansion behavior qualitatively matches Figs.~6--7, and the charge-2
$\to$ two-$(+1)$-core resolution with conserved circulation is the correct
topological realization of the paper's stability claim.

\medskip
\noindent\textbf{Verdict: PARTIAL.} Coverage 7/10, agreement 8/10.
\begin{itemize}
  \item \textbf{Coverage note:} straight-vortex statics + subcritical NLS
        dynamics + pressure-driven core breathing + charge-2 topological
        resolution are captured. The paper's \emph{headline} 3D vortex-ring
        splitting (Figs.~8--9) was not attempted.
  \item \textbf{Agreement note:} the quantitative static comparisons match to
        $\sim$1--2\%; the dynamic behavior matches qualitatively.
\end{itemize}

\section{Honest gaps}
\begin{enumerate}
  \item \textbf{No macroscopic separation.} The two $+1$ cores form but stay
        grid-adjacent ($\sim$1 healing length) rather than separating
        macroscopically. This is \emph{expected}: the paper itself notes a
        straight single vortex remains radially symmetric; the macroscopic split
        is shown for vortex \emph{rings}, where velocity-field asymmetry drives
        it (Figs.~8--9), a 3D axisymmetric calculation not reimplemented here.
  \item \textbf{Approximate initial $s=2$ core.} The seed is not the exact $s=2$
        solution of Eq.~\eqref{eq:profile}, so a little breathing occurs even
        undriven.
  \item \textbf{No quantitative growth-rate comparison.} The paper reports the
        core-energy parameter $\ell(t)$ (Eq.~42) and ring counts, not a linear
        split growth rate, so no head-to-head rate is available.
\end{enumerate}

\section{Reproducibility}
All code, the result JSON, and this report are self-contained. Part~A:
\texttt{vortex\_profile.py}. Part~B: \texttt{vortex\_split\_2d.py}. Driver /
verdict emitter: \texttt{save\_result.py} (writes
\texttt{berloff2008\_result.json}). Python via numpy/scipy
(\texttt{/home/stevens/comfyui-env/bin/python}). See
\texttt{report/evidence/} for copies of all code and the result JSON, and
\texttt{report/workflow.md} for tools/versions and effort.

\end{document}
